\section{Reproduction and Compile-Time Contract}
\label{app:reproduction}

\subsection{HAO-grid suite}

Run from \code{tk\_fa4/fp4\_fa4\_fwd}. The following command shows one
representative shape and policy; repeat \code{--shape} for the rows in
Eq.~\ref{eq:hao-grid} and run \code{nvmx-fast} and
\code{nvmx-accurate} in separate output directories.
\begin{lstlisting}[language=bash]
CUDA_VISIBLE_DEVICES=0 python3 hao_comprehensive_suite.py \
  --shape 1,4096,24,128 \
  --variant nvmx-fast \
  --warmup-ms 300 --rep-ms 3000 \
  --cooldown-seconds 0.8 \
  --seed 20260814 --gpu 0 \
  --output-dir /path/to/fast_shard0 \
  --build-root /tmp/fp4_fa4_fast_shard0
\end{lstlisting}
The production replay used four identical GB200s. Each shard contains its
own manifest; the generated summary validates exactly one complete case per
shape and policy:
\begin{lstlisting}[language=bash]
python3 \
  results/fp4_fa4_hao_table_gb200_20260802/build_summary.py
\end{lstlisting}

\subsection{FP8 and saturation controls}

The independent wider suite adds FP8 PV and H64 shapes:
\begin{lstlisting}[language=bash]
python3 hao_comprehensive_suite.py \
  --variant all \
  --shape 1,256,16,128 \
  --shape 1,1024,24,128 \
  --shape 1,2048,24,128 \
  --shape 1,4096,24,128 \
  --shape 1,4096,64,128 \
  --shape 1,8192,64,128 \
  --warmup-ms 300 --rep-ms 3000 \
  --seed 20260814 --gpu 0 \
  --output-dir /path/to/fp4_fa4_unified_replay
\end{lstlisting}
Its canonical BF16/HAO reference harness preallocates outputs and cycles
through all six provider orders. Every order and per-window median remains
in the JSON audit record.

\subsection{Downstream provider suite}

The downstream matrix uses one adapter for TK and native HAO. Every provider
receives identical weights, dataset order, seed, padding mask, and prepared
Q/K/V values.
\begin{lstlisting}[language=bash]
cd tk_fa4/fp4_fa4_fwd
python3 downstream_provider_suite.py \
  --gpu 0 \
  --output-dir \
  ../../results/fp4_fa4_downstream_matrix_20260801
\end{lstlisting}
Tasks and providers can be sharded with repeated \code{--task} and
\code{--provider}. Regenerate the audit without rerunning models using
\code{--summarize-only}. The suite hashes every full BF16 replay and checks
a matching BF16 prefix for the fail-fast NV/NV control. HAO uses its native
CuTe forward, not the TK port.

\subsection{Wan QK-guarded policy bundle}

Wan uses a layer-routed bundle rather than one global fallback. The build
script compiles the fast and accurate NV/MX bases, compiles the wide QK guard,
and records the routing in one manifest. The following replays the 14B fast
policy without exact stable softmax:
\begin{lstlisting}[language=bash]
cd tk_fa4/fp4_fa4_fwd
python3 build_wan_nv_mx_bundle.py \
  --model 14b \
  --output-dir /tmp/wan_nv_mx_14b_s7680
CUDA_VISIBLE_DEVICES=0 \
PYTHONPATH=/workspace/codebases/flash-attention-fp4 \
python3 eval_wan_video.py \
  --model Wan-AI/Wan2.1-T2V-14B-Diffusers \
  --providers hao-bf16,tk \
  --policy-manifest \
    /tmp/wan_nv_mx_14b_s7680/manifest.json \
  --policy fast --steps 4 \
  --output /tmp/wan_14b_fast.json
\end{lstlisting}
Use \code{--policy accurate} for the represented-normalization base. The
manifest selects guarded layers 27--29 for 1.3B and 33--34/38--39 for 14B;
all other layers remain on the selected base policy. The default guard uses
margin/stored-shift pairs $110/16$ and $112/14$, respectively. The builder
rejects a combined shift above 126, and the kernel clamps translated E8M0
scales to code 1 before evaluating represented normalization in the
non-underflowing order.

The main Wan table and its CuTe-BF16 paired outputs are in
\path{results/fp4_fa4_wan_cute_bf16_20260806}. Rebuild the table with:
\begin{lstlisting}[language=bash]
python3 results/fp4_fa4_wan_cute_bf16_20260806/\
build_tables.py
\end{lstlisting}
The corrected guard manifests, paired prompt replays, and five-second timing
windows are in
\path{results/fp4_fa4_wan_guard_fix_20260806}.

The layer-boundary calibration artifacts and aggregation script are in
\path{results/fp4_fa4_wan_affine_calibration_20260806}. The route evaluator's
\code{--active ROUTE=LAYERS} option runs FP4 only in selected layers and BF16
attention elsewhere, which keeps every candidate on the same teacher
trajectory. Rebuild the audit table with:
\begin{lstlisting}[language=bash]
python3 results/fp4_fa4_wan_affine_calibration_20260806/\
build_summary.py
\end{lstlisting}

The rejected end-to-end joint route search, wider prompt validation, and
fresh-process controls are in
\path{results/fp4_fa4_wan_joint_20260806}. The corresponding reusable tools
are \path{joint_optimize_wan_affine.py},
\path{validate_wan_joint_routes.py}, and
\path{materialize_wan_joint_route.py}.

\subsection{ViT-MAE reconstruction}

Download the 100 recorded COCO 2017 validation filenames into one directory,
then run each provider in a separate process so its extension and CuTe cache
remain isolated. A representative fast run is:
\begin{lstlisting}[language=bash]
cd tk_fa4/fp4_fa4_fwd
EXT=/tmp/tk_hao_unified_g0/
EXT=${EXT}b1_s256_h16_d128_nvmx-fast.so
CUDA_VISIBLE_DEVICES=0 python3 \
  eval_vit_mae_reconstruction.py \
  --image-dir /path/to/coco-val-images \
  --samples 100 \
  --extension $EXT \
  --output /tmp/nvmx_fast_100.json
\end{lstlisting}
For the accurate route, add \code{--global-anchor-kv}. For HAO controls, set
\code{--attention-backend} to \code{hao-native} or \code{hao-fp8}. The input
order is stored in every JSON file. Regenerate paired intervals and the table
with:
\begin{lstlisting}[language=bash]
python3 results/fp4_fa4_reconstruction_20260805/\
build_summary.py
\end{lstlisting}

\subsection{B300 mini-job}

The committed compatibility job can be submitted with the Volt production
CLI after exposing the existing \code{GITHUB\_TOKEN} secret:
\begin{lstlisting}[language=bash]
volt job submit \
  results/fp4_fa4_b300_tuning_20260802/volt_b300_smoke.yaml \
  --tag project:fp4-fa4 --tag purpose:report
\end{lstlisting}
The job clones fixed TK and HAO commits, builds for \code{sm\_103a}, and
stores hardware metadata, logs, the extension hash, and benchmark JSON under
Volt artifacts.

The follow-up job corrects the persistent-grid width and compiles four
native-EX2 densities from the same source and input tensors:
\begin{lstlisting}[language=bash]
volt job submit \
  results/fp4_fa4_b300_tuning_20260802/\
volt_b300_density_sweep.yaml \
  --tag project:fp4-fa4 --tag purpose:report
\end{lstlisting}

The complete D64 runs and final no-override policy check use:
\begin{lstlisting}[language=bash]
volt job submit results/fp4_fa4_b300_tuning_20260802/\
volt_b300_d64_nvmx_full_sweep.yaml
volt job submit results/fp4_fa4_b300_tuning_20260802/\
volt_b300_d64_nvnv_full_sweep.yaml
volt job submit results/fp4_fa4_b300_tuning_20260802/\
volt_b300_d64_promoted_policy.yaml
\end{lstlisting}
The corresponding successful job IDs are \path{f8ajtqmbvd8p},
\path{ig2tvyqoxtlz}, and \path{2bz33uyyvvo8}. The grid-cap sweep is
\path{q6nb0jj73xik}; bounded NV/NV validation is \path{szpe9052wipe}.
The final D128 cross-generation rows come from the order/KV-stage replay
\path{ev0g6vopwujl}, saturated SFU/reduction sweep \path{y5yrv3k9dvxh},
S6144 prefetch/SFU sweep \path{oen3rbbipbtf}, and long-context job
\path{xfby78ic4nql}. The 3-PFLOP/s discovery sweep is
\path{ibwp4shzwfmi}; its independent five-window, two-seed confirmation is
\path{d9e3yhudpu93}. Their specifications are
\path{volt_b300_d128_3ktflops_sweep.yaml} and
\path{volt_b300_d128_3ktflops_confirm.yaml}.
After downloading the artifacts, regenerate the B300 JSON and LaTeX rows by
running \code{python3 build\_summary.py} from the
\path{results/fp4_fa4_b300_tuning_20260802} directory.

\subsection{Required operand contract}

Both retained policies require folded K64 Q/K scales selected by the MSE
policy:
\begin{lstlisting}[language=bash]
--nv-qk-fold-k64-scales both \
--nv-qk-fold-scale-select mse
\end{lstlisting}
\code{Accurate} also requires the same fixed 32-row permutation for K and V.
Pairing a folded-scale binary with ordinary block-16 Q/K scales can produce
plausible timing and invalid accuracy; the suite couples the binary and
operand arguments.

\subsection{Retained compile-time policy}

\code{Fast} uses a 12-stage K/V ring, all-affine mode 23, four-word
represented denominators, paired scale reuse, wide three-input maxima,
Q-scale preloading, early asynchronous P-scale handoff, and a 200/208/56
register split. \code{Accurate} uses 13 K/V stages, native EX2 samples, a
32-row anchor, and correction-warpgroup normalization.

Historical NV/NV, intermediate NV/MX, sparse, sampled-denominator,
direct-code, quadratic, two-CTA, alternate-owner, and interleaving probes
remain compile-gated and default-off.

\section{Terminology}

\begin{description}
  \item[Quarter] One N32 fragment read from an N128 score tile. Four producer
  quarters feed two K64 scaled-FP4 PV commands.
  \item[Score bank] A 128-column FP32 TMEM region receiving QK and later
  hosting transient P/scale overlays.
  \item[Output bank] A permanent 128-column FP32 PV accumulator for one query
  stage.
  \item[Represented denominator] A normalizer summed from the E2M1 payload
  and decoded block scale actually consumed by PV.
  \item[Early publication] Signaling a legal first K64 payload and scale pair
  before tail P construction completes.
  \item[Speed-of-light diagnostic] An intentionally incorrect or
  semantically altered kernel that removes work to bound latency.
  \item[Full FP4] Q, K, P, and V use E2M1 payloads with block scales. NVFP4
  QK plus FP8 PV is a control, not full FP4.
\end{description}
